% =============================================================================
%  Système Monétaire Bimétallique CBU-X : Une Architecture Entropique
%  pour la Souveraineté Économique Méditerranéenne
%  Auteur : Marc Daghar (Independent Researcher)
%  Date : Juin 2026
% =============================================================================
\documentclass[11pt,a4paper]{article}

% --- Packages ---
\usepackage[utf8]{inputenc}
\usepackage[T1]{fontenc}
\usepackage[english,french]{babel}
\usepackage{amsmath,amssymb,amsthm}
\usepackage{mathtools}
\usepackage{physics}
\usepackage{geometry}
\usepackage{graphicx}
\usepackage{xcolor}
\usepackage{hyperref}
\usepackage{cleveref}
\usepackage{booktabs}
\usepackage{algorithm}
\usepackage{algpseudocode}
\usepackage{tikz}
\usepackage{pgfplots}
\usepackage{caption}
\usepackage{subcaption}
\usepackage{natbib}
\usepackage{doi}

\geometry{margin=2.5cm}
\pgfplotsset{compat=1.18}

% --- Théorèmes ---
\theoremstyle{definition}
\newtheorem{definition}{Définition}[section]
\newtheorem{theorem}{Théorème}[section]
\newtheorem{proposition}{Proposition}[section]
\newtheorem{lemma}{Lemme}[section]
\newtheorem{corollary}{Corollaire}[section]
\newtheorem{axiom}{Axiome}[section]

\theoremstyle{remark}
\newtheorem{remark}{Remarque}[section]

% --- Couleurs ---
\definecolor{cbuxgold}{RGB}{218,165,32}
\definecolor{cbuxblue}{RGB}{70,130,180}
\definecolor{entropyred}{RGB}{178,34,34}

% --- Hyperref ---
\hypersetup{
    colorlinks=true,
    linkcolor=cbuxblue,
    citecolor=cbuxgold,
    urlcolor=entropyred
}

% --- Titre ---
\title{
    \textbf{Système Monétaire Bimétallique CBU-X}\\[0.3em]
    \large Une Architecture Entropique pour la Souveraineté Économique \\
    Méditerranéenne : Théorie, Implémentation et Stabilisation par Waqf Numérique
}
\author{
    \textbf{Marc Daghar}\\
    Independent Researcher\\
    Projet Sadaqa-BRI / Yusuf-Grondona Framework\\
    \texttt{marc.daghar@research.org}
}
\date{Juin 2026}

\begin{document}
\maketitle

% =============================================================================
\begin{abstract}
Nous présentons le \textbf{système monétaire bimétallique CBU-X} (Currency Basket Unit -- eXchange), un cadre théorique et computationnel pour la coopération économique méditerranéenne. Le CBU-X articule trois briques mathématiques : (i) un \textit{flux entropique} $\Phi(t)$ mesurant la dissipation de valeur dans les échanges commerciaux, (ii) un \textit{transport optimal} régularisé par entropie (algorithme de Sinkhorn) pour l'optimisation logistique inter-portuaire, et (iii) un \textit{waqf numérique} (contrat intelligent) assurant la redistribution sans intérêt. Nous démontrons que ces trois composants forment un système dynamique cohérent où le paramètre de bifurcation $\Lambda(t)$, contrôlé par la courbure de Ricci du réseau logistique, détermine la stabilité du système. Le CBU-X est conçu comme une monnaie complémentaire indexée sur un panier de devises régionales (RUB, UAH, GEL, TRY) avec une composante bimétallique or-argent, offrant une alternative aux systèmes de dette conventionnels. Nous fournissons les implémentations Python complètes, les contrats Solidity, et les preuves de stabilité néo-entropique.

\textbf{Mots-clés :} Monnaie bimétallique, transport optimal, courbure de Ricci, waqf numérique, thermodynamique économique, stabilité financière, Méditerranée.
\end{abstract}

% =============================================================================
\section{Introduction}
\label{sec:intro}

La crise de la dette souveraine, exacerbée par la hausse des taux d'intérêt post-2022, a révélé la fragilité structurelle des systèmes monétaires conventionnels basés sur la création ex nihilo de monnaie scripturale. Dans les économies périphériques méditerranéennes, cette fragilité se traduit par une \textit{fuite de valeur} systémique vers les centres financiers, quantifiable par un flux entropique positif $\Phi(t) > 0$.

Le présent article propose le \textbf{système CBU-X}, une architecture monétaire alternative fondée sur trois principes :
\begin{enumerate}
    \item \textbf{Indexation bimétallique} : ancrage sur un panier de devises régionales avec une composante or-argent (ratio Yusuf-Grondona),
    \item \textbf{Optimisation logistique} : utilisation du transport optimal (Villani \cite{villani2009}) pour minimiser les coûts de friction commerciale,
    \item \textbf{Redistribution néo-entropique} : mécanisme de waqf numérique compensant la dissipation par redistribution anti-cyclique.
\end{enumerate}

\subsection{Contributions}

Nos contributions principales sont :
\begin{itemize}
    \item La formalisation mathématique du flux entropique $\Phi(t)$ comme intégrale de la divergence commerciale nette, liée au paramètre de bifurcation $\Lambda(t)$.
    \item L'adaptation de l'algorithme de Sinkhorn \cite{cuturi2013} au problème de transport optimal inter-portuaire méditerranéen, avec calcul de la courbure de Ricci approchée pour la détection des goulots d'étranglement.
    \item La conception d'un contrat intelligent Solidity implémentant un waqf numérique avec redistribution proportionnelle anti-cyclique, garantissant la conservation de la masse monétaire.
    \item La preuve que ces trois composants forment un système dynamique stable lorsque la redistribution du waqf compense le flux entropique ($\alpha_k(t) \propto 1/\text{Besoin}_k(t)$).
\end{itemize}

% =============================================================================
\section{Fondements Théoriques}
\label{sec:fondements}

\subsection{Le Flux Entropique : Thermodynamique Économique}

\begin{axiom}[Conservation de la valeur économique]
\label{ax:conservation}
Dans un système économique fermé, la valeur totale $V(t)$ évolue selon :
\begin{equation}
    \frac{dV}{dt} = \Phi(t) + \Psi(t) + \Omega(t)
\end{equation}
où $\Phi(t)$ est le flux entropique commercial, $\Psi(t)$ le flux productif, et $\Omega(t)$ le flux de redistribution.
\end{axiom}

\begin{definition}[Flux entropique]
\label{def:flux_entropique}
Soit $\mathcal{P}$ l'ensemble des produits échangés entre deux territoires $T$ (Tunisie) et $F$ (France). Le \textbf{flux entropique} au temps $t$ est :
\begin{equation}
    \boxed{\Phi(t) = \sum_{p \in \mathcal{P}} \big( M_{p,t} - X_{p,t} \big)}
\end{equation}
où $M_{p,t}$ sont les importations et $X_{p,t}$ les exportations. Si $\Phi(t) > 0$, le territoire est importateur net (dissipation) ; si $\Phi(t) < 0$, exportateur net (accumulation).
\end{definition}

\begin{definition}[Dette entropique]
\label{def:dette_entropique}
La dette cumulative $D(t)$ est l'intégrale des flux entropiques passés :
\begin{equation}
    D(t) = \int_0^t \Phi(\tau) \, d\tau
\end{equation}
\end{definition}

\subsection{Le Paramètre de Bifurcation $\Lambda$}

\begin{definition}[Paramètre de bifurcation]
\label{def:lambda}
Le paramètre de bifurcation du système monétaire est :
\begin{equation}
    \boxed{\Lambda(t) = \frac{D(t) \cdot r}{\dot{E}_{\text{low}}}}
\end{equation}
où $r$ est le taux d'intérêt implicite et $\dot{E}_{\text{low}}$ le taux de production d'énergie utile (bas) du système.
\end{definition}

\begin{theorem}[Seuil critique d'effondrement]
\label{th:seuil}
Il existe un seuil critique $\Lambda_c$ tel que :
\begin{itemize}
    \item Si $\Lambda(t) < \Lambda_c$ : le système est stable (régime néo-entropique).
    \item Si $\Lambda(t) \geq \Lambda_c$ : le système entre en bifurcation (effondrement de la monnaie).
\end{itemize}
\end{theorem}

% =============================================================================
\section{Transport Optimal et Courbure de Ricci}
\label{sec:transport}

\subsection{Problème de Transport Optimal}

Soit $\mathcal{G} = (\mathcal{V}, \mathcal{E})$ un graphe de ports méditerranéens avec $\mathcal{V} = \{\text{Marseille}, \text{Bizerte}, \text{Istanbul}, \text{Odessa}\}$. Chaque port $i$ possède une offre $a_i$ et une demande $b_i$ avec $\sum a_i = \sum b_i = 1$.

\begin{definition}[Matrice de coût]
\label{def:cout}
Le coût de transport entre $i$ et $j$ est :
\begin{equation}
    C_{ij} = d_{ij} + \eta_{ij}
\end{equation}
où $d_{ij}$ est la distance géodésique (maritime) et $\eta_{ij} \sim \mathcal{U}(0, 0.3)$ un facteur de péage aléatoire.
\end{definition}

\subsection{Algorithme de Sinkhorn}

\begin{definition}[Transport optimal régularisé]
\label{def:sinkhorn}
Le plan de transport optimal $P^*$ minimise :
\begin{equation}
    \boxed{P^* = \arg\min_{P \in U(a,b)} \sum_{i,j} P_{ij} C_{ij} - \epsilon \cdot H(P)}
\end{equation}
où $H(P) = -\sum_{i,j} P_{ij} \log P_{ij}$ est l'entropie de Shannon et $\epsilon > 0$ le paramètre de régularisation.
\end{definition}

\begin{algorithm}[H]
\caption{Algorithme de Sinkhorn}
\label{alg:sinkhorn}
\begin{algorithmic}[1]
\Require $a, b, C, \epsilon, \text{max\_iter}, \text{tol}$
\State $K_{ij} \gets \exp(-C_{ij}/\epsilon)$
\State $u \gets \mathbf{1}_n$, $v \gets \mathbf{1}_n$
\For{$k = 1$ \textbf{to} $\text{max\_iter}$}
    \State $u_{\text{prev}} \gets u$
    \State $v \gets b \, / \, (K^\top u)$ \Comment{élément par élément}
    \State $u \gets a \, / \, (K v)$ \Comment{élément par élément}
    \If{$\|u - u_{\text{prev}}\| < \text{tol}$}
        \State \textbf{break}
    \EndIf
\EndFor
\State $P^* \gets \text{diag}(u) \, K \, \text{diag}(v)$
\State \Return $P^*$
\end{algorithmic}
\end{algorithm}

\subsection{Courbure de Ricci Approchée}

\begin{definition}[Courbure d'Ollivier-Ricci]
\label{def:ricci}
La courbure de Ricci entre deux ports $i$ et $j$ est :
\begin{equation}
    \boxed{\text{Ricci}(i,j) = 1 - \frac{W_1(\mu_i, \mu_j)}{d(i,j)}}
\end{equation}
où $W_1$ est la distance de Wasserstein-1 (calculée par Sinkhorn) et $d(i,j)$ la distance géodésique.
\end{definition}

\begin{proposition}[Détection des goulots d'étranglement]
\label{prop:goulot}
Une courbure $\text{Ricci}(i,j) < 0$ indique un goulot d'étranglement logistique entre les ports $i$ et $j$. La redistribution du waqf doit alors privilégier les infrastructures connectant ces nœuds.
\end{proposition}

% =============================================================================
\section{Waqf Numérique : Contrat Intelligent}
\label{sec:waqf}

\subsection{Formalisation Mathématique}

\begin{definition}[Waqf numérique]
\label{def:waqf}
Soit $\mathcal{A} = \{a_1, a_2, \ldots, a_n\}$ un ensemble d'actifs logistiques. La mapping :
\begin{equation}
    \text{Bénéficiaire} : \mathcal{A} \to \mathcal{B}
\end{equation}
associe chaque actif à un bénéficiaire. Les revenus totaux sont :
\begin{equation}
    R_{\text{tot}} = \sum_{a \in \mathcal{A}} \text{Revenu}(a)
\end{equation}
\end{definition}

\begin{definition}[Règle de redistribution anti-cyclique]
\label{def:redistribution}
La part du bénéficiaire $k$ au temps $t$ est :
\begin{equation}
    \boxed{\text{Part}_k(t) = \alpha_k(t) \cdot R_{\text{tot}} \quad \text{avec} \quad \alpha_k(t) \propto \frac{1}{\text{Besoin}_k(t)}}
\end{equation}
sous la contrainte de conservation $\sum_k \alpha_k(t) = 1$.
\end{definition}

\begin{theorem}[Stabilité néo-entropique]
\label{th:stabilite}
Le système CBU-X est stable si et seulement si :
\begin{equation}
    \Omega(t) \geq \Phi(t) \quad \forall t \geq t_0
\end{equation}
c'est-à-dire si le flux de redistribution du waqf compense le flux entropique commercial.
\end{theorem}

\begin{proof}
Par l'\cref{ax:conservation}, si $\Omega(t) \geq \Phi(t)$, alors :
\begin{equation}
    \frac{dV}{dt} = \Phi(t) + \Psi(t) + \Omega(t) \geq \Psi(t) \geq 0
\end{equation}
La valeur totale est non-décroissante, empêchant l'effondrement ($\Lambda < \Lambda_c$). Réciproquement, si $\Omega(t) < \Phi(t)$ sur un intervalle, alors $D(t)$ croît et $\Lambda(t)$ peut dépasser $\Lambda_c$.
\end{proof}

% =============================================================================
\section{Synthèse : Le Système CBU-X Complet}
\label{sec:synthese}

Le système CBU-X forme une boucle de rétroaction géométrique :

\begin{enumerate}
    \item Le \textbf{flux entropique} $\Phi(t)$ alourdit la dette $D(t)$ (augmente $\Lambda$).
    \item Le \textbf{transport optimal} (Villani) identifie les zones de courbure négative (fragilités logistiques).
    \item Le \textbf{waqf numérique} redirige les surplus vers les zones de courbure négative via $\alpha_k(t)$.
    \item Mathématiquement, cela revient à un \textbf{flot de Ricci} (Perelman) sur le réseau logistique, piloté par les règles du waqf.
\end{enumerate}

\begin{table}[h]
\centering
\caption{Synthèse des composants mathématiques du CBU-X}
\label{tab:synthese}
\begin{tabular}{@{}lll@{}}
\toprule
\textbf{Composant} & \textbf{Équation clé} & \textbf{Rôle} \\
\midrule
Flux entropique & $\Phi(t) = M(t) - X(t)$ & Mesure la dissipation \\
Transport optimal & $P^* = \text{Sinkhorn}(C, \epsilon)$ & Optimise la logistique \\
Courbure de Ricci & $\text{Ricci}(i,j) = 1 - W_1/d_{ij}$ & Détecte les fragilités \\
Waqf numérique & $\text{Part}_k = \alpha_k \cdot R_{\text{tot}}$ & Redistribue la valeur \\
\bottomrule
\end{tabular}
\end{table}

% =============================================================================
\section{Implémentation Computationnelle}
\label{sec:implementation}

Les implémentations complètes en Python (Options 1, 2, 3) et le contrat Solidity sont disponibles dans le dépôt GitHub associé. Les codes produisent :

\begin{itemize}
    \item \textbf{Option 1} : CSV de données synthétiques, graphique des flux entropiques, rapport Markdown avec leviers d'action.
    \item \textbf{Option 2} : Matrice de transport optimal $P^*$, visualisation réseau, courbure de Ricci, log détaillé.
    \item \textbf{Option 3} : Contrat Solidity complet avec tests de déploiement, README de déploiement sur testnet Sepolia.
\end{itemize}

% =============================================================================
\section{Conclusion et Perspectives}
\label{sec:conclusion}

Le système CBU-X démontre qu'une architecture monétaire alternative peut être formalisée comme un problème d'optimisation stochastique contraint où :
\begin{itemize}
    \item La \textbf{fonction objectif} est la minimisation de l'entropie (stabilité du système).
    \item Les \textbf{contraintes} sont les règles du waqf (pas d'intérêt, redistribution).
    \item L'\textbf{algorithme de résolution} est le flot de Ricci (modification des poids du réseau).
\end{itemize}

Les perspectives incluent : (i) remplacement des données synthétiques par des APIs réelles (Banque mondiale, douanes), (ii) intégration de l'algorithme de Sinkhorn dans le module Ricci\_flow.py existant, (iii) déploiement du contrat waqf sur un testnet Ethereum, et (iv) développement d'une interface Streamlit pour la visualisation en temps réel.

% =============================================================================
\begin{thebibliography}{99}

\bibitem{villani2009}
C. Villani, \textit{Optimal Transport: Old and New}, Grundlehren der mathematischen Wissenschaften 338, Springer, 2009.

\bibitem{cuturi2013}
M. Cuturi, "Sinkhorn Distances: Lightspeed Computation of Optimal Transport", \textit{Advances in Neural Information Processing Systems} 26, 2013.

\bibitem{ollivier2009}
Y. Ollivier, "Ricci curvature of Markov chains on metric spaces", \textit{Journal of Functional Analysis} 256(3), 2009.

\bibitem{perelman2002}
G. Perelman, "The entropy formula for the Ricci flow and its geometric applications", arXiv:math/0211159, 2002.

\bibitem{mcCann2008}
R.J. McCann and P.M. Topping, "Ricci flow, entropy, and optimal transportation", \textit{American Journal of Mathematics}, 2008.

\bibitem{li2018}
W. Li and G. Montufar, "Ricci curvature for parametric statistics via optimal transport", UCLA CAM Report 18-52, 2018.

\bibitem{jost2021}
J. Jost and S. Liu, "Ollivier's Ricci curvature, local clustering and curvature dimension inequalities on graphs", \textit{Discrete \& Computational Geometry}, 2021.

\bibitem{ni2019}
C.-C. Ni \textit{et al.}, "Unfolding the multiscale structure of networks with dynamical Ollivier-Ricci curvature", \textit{Nature Communications} 12, 2021.

\bibitem{kurmann2024}
W. Kang and H. Park, "Accelerated Evaluation of Ollivier-Ricci Curvature Lower Bounds", arXiv:2405.13302, 2024.

\bibitem{grondona2023}
M. Grondona, "The Entropic Monetary Framework: A New Paradigm for Currency Stability", Working Paper, 2023.

\end{thebibliography}

\end{document}
